% ============================================================================
%  D — file phần: Sắc thái
%  Ngày tạo: 2026-09-07 | Sửa gần nhất: 2026-09-07 | Ôn gần nhất: chưa
% ============================================================================
\documentclass[../english.tex]{subfiles}
\begin{document}
\part{SẮC THÁI}
\begin{tomtat}
Phần D đi vào đúng mục tiêu bạn đặt ra ban đầu: hiểu sâu về ngôn từ. Ba chương trả lời ba câu hỏi mà từ điển không trả lời được. Khi hai từ ``cùng nghĩa'', cái gì quyết định chọn cái nào (16)? Vì sao một số tổ hợp từ nghe tự nhiên và số khác thì không, dù không luật nào cấm (17)? Và vì sao cách một ngành nói về đối tượng của nó lại định hình cách người trong ngành nghĩ (18)? Nếu chỉ nhớ một điều: nghĩa của một từ không nằm trong định nghĩa của nó mà nằm trong \emph{các lựa chọn mà người nói đã bỏ qua} để chọn nó.
\end{tomtat}

Phần A đã đặt nền: nghĩa nằm trong cách dùng, và biết một từ gồm chín mặt chứ không phải một. Phần B và C dùng nền đó cho việc đọc và viết. Phần này quay lại chính cái nền và đào sâu ba lớp mà người học tiếng Anh thường dừng lại trước.

Lớp thứ nhất là \emph{lựa chọn giữa các từ gần nghĩa}. Đây là chỗ từ điển song ngữ phản bội bạn nhiều nhất: nó cho ba từ tiếng Anh cho một từ tiếng Việt và không nói gì về việc chọn cái nào. Chương 16 cho hai công cụ tách chúng --- văn vực và các chiều sắc thái --- và một cách tra cứu thực tế bằng dữ liệu ngôn ngữ thật.

Lớp thứ hai là \emph{kết hợp từ}. Đây là thứ phân biệt rõ nhất người viết thành thạo với người viết đúng ngữ pháp. \emph{Do a decision} không sai luật nào cả; nó chỉ không phải cách người ta nói. Chương 17 giải thích vì sao hiện tượng này tồn tại, vì sao nó quan trọng hơn ngữ pháp với người đọc bản ngữ, và cách học nó mà không phải thuộc lòng.

Lớp thứ ba là \emph{ẩn dụ}. Nghe như chuyện văn chương, nhưng nó là chuyện kỹ thuật: ngành nào cũng nói về đối tượng của mình qua một hệ ẩn dụ, và hệ đó vừa cho phép vừa giới hạn cách nghĩ. Chương 18 chỉ ra chúng trong chính ngành của bạn.
\subfile{00-map}
\subfile{16-gan-nghia-va-van-vuc/gan-nghia-va-van-vuc}
\subfile{17-ket-hop-tu/ket-hop-tu}
\subfile{18-an-du-va-khung-tu-duy/an-du-va-khung-tu-duy}
\ifSubfilesClassLoaded{\printbibliography[title={Nguồn}]}{}
\end{document}
